\documentclass[11pt]{article}

\usepackage[margin=1in]{geometry}
\usepackage{enumitem}
\usepackage{xcolor}
\usepackage{mdframed}
\usepackage{parskip}
\usepackage{hyperref}
\hypersetup{colorlinks=true, linkcolor=black, urlcolor=black, citecolor=black}

% Environment for a committee comment
\newmdenv[
  linecolor=black!25,
  linewidth=0.6pt,
  backgroundcolor=black!4,
  skipabove=10pt,
  skipbelow=6pt,
  innertopmargin=6pt,
  innerbottommargin=6pt,
  leftmargin=0pt,
  rightmargin=0pt,
]{commentbox}

\newcounter{commentcounter}
\newcommand{\comment}[1]{%
  \refstepcounter{commentcounter}%
  \begin{commentbox}
  \textbf{Comment \thecommentcounter.}\quad\itshape #1
  \end{commentbox}
}

\setlist[itemize]{leftmargin=1.4em, itemsep=3pt, topsep=4pt}

\title{\vspace{-2em}Addressing the Committee's Comments}
\author{Zhiyang Chen}
\date{\today}

\begin{document}
\maketitle

This letter summarizes the revisions made to Chapter~1 (Introduction) in
response to the committee's comments. Page and section numbers refer to the
current draft (\texttt{main.pdf}).

%========================================================
\comment{We do not need current Section~1.1. The shared umbrella is not good
enough. It reads like parallel stories forced into one.}

\begin{itemize}
  \item I removed the original preamble. Section~1.1 is now a Motivation
  section (pages 2--5). Instead of forcing smart contracts and LLMs into one
  abstract umbrella, it follows a single software lifecycle (Figure~1.1): the
  same software is first built, now with AI assistance, and then deployed, now
  used by autonomous agents. The two lines of work (LLM security and smart
  contract security) are no longer parallel stories; they are separate stages
  of one lifecycle, which I call \emph{untrusted development} (Section~1.1.2)
  and \emph{untrusted use} (Section~1.1.3). The section opens from concrete
  examples rather than an abstract framing.

  \item I kept the detailed related work spread across the individual chapters,
  where each technical chapter still has its own related-work section. In
  addition, Chapter~1 now includes a short Prior Work section (Section~1.3)
  that gives a high-level overview of what prior work looked like before this
  dissertation. This overview lets an outside reader see the landscape and
  where this dissertation fits.
\end{itemize}

%========================================================
\comment{When we talk about emerging software, it is obviously tied to the
blockchain side. There is a series of blockchain-poisoning incidents that can
serve as part of the motivation.}

\begin{itemize}
  \item I replaced ``Emerging Software'' with ``High-Stakes Software'', and
  retitled the thesis to \emph{Securing High-Stakes Software Against Untrusted
  Development and Untrusted Use}. Section~1.1.1 now states clearly that
  blockchain applications are the representative example of high-stakes
  software, since they hold stakes (cryptocurrencies) directly in code and a
  single compromise is an immediate, measurable loss (that is why I chose it as
  the representative).

  \item I also added a series of cybersecurity incidents to the motivation
  itself. All three are real, high-financial-loss cybersecurity incidents, and
  each one comes from a mistake at a different stage of the development
  lifecycle:
  \begin{enumerate}[leftmargin=1.8em, itemsep=2pt, topsep=3pt]
    \item at the coding stage, an AI assistant wrote a scam API endpoint into a
    code snippet and the developer lost about \$2{,}500 (Section~1.2.1);
    \item an AI-co-authored pull request introduced a faulty price oracle that
    passed AI review and cost the Moonwell protocol about \$1.78M
    (Section~1.2.2); and
    \item at the deployment (use) stage, carefully written and audited code was
    subverted through an unanticipated interaction, costing Harvest Finance
    about \$33.8M (Section~1.2.3).
  \end{enumerate}
  Together they show that a single mistake at any stage of the lifecycle can
  lead directly to a large, measurable financial loss.
\end{itemize}

%========================================================
\comment{There are too many terminologies in the first chapters, which makes
them annoying to read. Imagine the audience as not only the committee but also
the public: one goal of a thesis is to enable outsiders. There is one step of
difference between the paper audience and the thesis audience.}

\begin{itemize}
  \item I rewrote the opening for a general reader and removed the paper-level
  jargon (for example ``Layer~I to~IV'', ``RAG-enabled coding agents'', and
  ``invariant guards''). The whole introduction now rests on two plain ideas,
  defined in everyday words in Section~1.1: \emph{development is untrusted} when
  the things used to build the software cannot be assumed safe, standing for
  AI-assisted coding tools (Section~1.1.2); \emph{use is untrusted} when the
  deployed software can be invoked by anyone unknown at deployment time, and
  the case studied in this thesis is smart contracts that anyone on the
  blockchain can call (Section~1.1.3).
\end{itemize}

%========================================================
\comment{The motivation does not need to align with the scope. It should be
more self-contained.}

\begin{itemize}
  \item I made the motivation self-contained and separated it from the scope.
  Section~1.2 now presents the three incidents as standalone stories that stand
  on their own evidence, and each ends in a general finding rather than a
  mapping to a specific chapter. The narrowing to this thesis's own problems
  and systems happens only afterward, in Section~1.2.4 (Problems This
  Dissertation Addresses) and Section~1.4 (Contributions). I also added a clear
  statement at the end of Section~1.2.4 that these questions and defenses are
  not tied to any single domain, and that blockchain is only the representative
  grounding. The motivation therefore reads as a general argument, not a
  restatement of the scope.

  \item To connect the general motivation back to this thesis, the
  contributions figure (Figure~1.7 in Section~1.4) reuses the same lifecycle
  diagram from Figure~1.1 and simply annotates it with this dissertation's
  contributions, one at each stage where risk is introduced.
\end{itemize}

\end{document}
